\anexo{Transcripción de la entrevista a Pablo Rodolfo Villarino}

\textbf{Entrevistado:} Pablo (ingeniero DevOps, IBM Argentina)

\textbf{Entrevistadores:} Muntaabski Federico, Mociulsky Santiago Bernardo

\begin{description}[leftmargin=0cm, labelsep=0.5cm]
    \item[\textbf{MUNTAABSKI FEDERICO:}] Bueno.

    \item[\textbf{PABLO:}] Perfect.

    \item[\textbf{MOCIULSKY SANTIAGO BERNARDO:}] Genial, ahí está.

    \item[\textbf{MUNTAABSKI FEDERICO:}] Perfecto, bien, bueno, ahora vamos a arrancar con la entrevista al profe Pablo. Vamos a arrancar con lo que sería el primer punto, que sería ¿Cómo describirías el mayor desafío de seguridad que hoy vos ves en la actualidad o que tuviste en tus organizaciones o empresas en las que trabajaste?

    \item[\textbf{PABLO:}] Perfecto. Bueno, primero contarles que actualmente me desempeño como ingeniero DevOps en IBM Argentina para los clientes de IBM Argentina. Yo trabajo aquí hace más o menos unos 14 años y me he movido por diferentes roles dentro de la compañía. Puntualmente, este rol dentro del mundo DevOps, el ingeniero DevOps se caracteriza por ser un rol donde nosotros tenemos que integrar un montón de software para armar plataformas de automatización que le permitan a los proyectos de software poder. Bueno, , como ustedes sabrán, poder confeccionar diferentes partes de aplicaciones o software, testear esos códigos, promulgarlos a otros ambientes de ejecución y todo eso habitualmente no se hace con un único software, sino que se integran herramientas. Tomando un poco tu pregunta, Federico, yo creo, y desde el lugar donde estoy trabajando actualmente, creo que uno de los puntos más o los desafíos más importantes en la seguridad es que justamente esa integración de herramientas diferentes que tienen y siguen estándares. Diferentes de desarrollo pueda ser realmente segura. Lo que lo que ocurre habitualmente es que podemos integrar. Podemos hacer que esas herramientas se vean entre sí, pero muchas veces esas herramientas son on premise, están dentro de nuestra red. o a veces dentro de la red de los clientes o muchas veces son, vamos a decir, plataformas online, es decir, las plataformas cloud como pueden ser GitLab o GitHub y ahí empiezan necesariamente que aparecer métodos en los cuales vamos a tener que integrar esas herramientas Pero a la vez, segurizar las comunicaciones entre las diferentes piezas de la misma plataforma, ¿no? Y entonces eso es un desafío porque. No todas las herramientas son compatibles con todas las técnicas de segurización de datos y muchas veces hay que hacer un poco de tuneo para que esa integración realmente sea segura en este tipo de proyectos DevOps y DevSecOps. El problema es que. Si nosotros no segurizamos correctamente ese ensamblaje de herramientas, esa integración, bueno, empiezan a aparecer superficies de ataque en cada unión. Es como si tuvieras una sábana toda llena de parches por todos lados y eso no está bien cocido. El día que en esta parte fuerte rompiste todo, ¿no? Entonces acá pasa un poco eso, hay muchas integraciones, de acuerdo al proyecto podemos tener 20, 25 herramientas diferentes y si eso no está bien integrado y bien segurizado, es un lío. Así que en mi expertise creo que ese es uno de los grandes desafíos, el primero por lo menos.

    \item[\textbf{MUNTAABSKI FEDERICO:}] Valora. Es Sandy.

    \item[\textbf{MOCIULSKY SANTIAGO BERNARDO:}] Sí, bueno, yo te quería consultar también si en uno de estos problemas ves involucrado la actividad del usuario. O sea, si como repregunta lo primero, si alguno de estos principales problemas es alguna interacción o algún factor en el que puede intervenir los usuarios que usan el software en el día a día.

    \item[\textbf{PABLO:}] Bueno, acá en mi especialidad en particular, quien va a usar esta plataforma para poder desarrollar aplicaciones, testearlas, etcétera. Básicamente son usuarios con expertise técnico informático, es decir, no es el típico usuario final el que va a usar. El software completando ventanas o campos, sino que son otros ingenieros, licenciados o licenciadas que conocen del mundo IT que se encargan de una parte del ciclo de desarrollo de software. Entonces, no sé, por ejemplo, los ingenieros DevOps nos encargamos de hacer todo ese ensamblaje de herramientas de software. y de mantenerlas, pero después van a venir los testerqua y solo van a utilizar las herramientas de testeo o van a venir los especialistas en observabilidad y monitoreo que solo van a utilizar las herramientas de observabilidad y monitoreo. Entonces, si tomo tu pregunta, Santiago, para mí usuarios es esa persona que va a usar una parte de esta plataforma de integración Divos y es ese perfil, esa persona, es una persona que sabe del mundo de tecnología, ¿no? Entonces, si la integración. Y las herramientas que usamos como integración. ¿Están correctamente configuradas? Cada una de estas herramientas permite lo que se llama segregación de usuarios, es decir, cada una de las herramientas le va a otorgar un rol a cada usuario y ese rol solo va a tener los permisos únicos y necesarios Para las acciones que debe hacer ese usuario, no sobre todo en plataformas cloud, los usuarios dentro de un mismo proyecto van a tener un usuario y una password que solo les va a permitir y autorizar a hacer una limitada cantidad de acciones, ¿no? Y esa, como ustedes sabrán, también esas acciones van a ser. trazables dentro de la plataforma, es decir, la misma plataforma, en este caso GitHub, que es una parte de esta gran plataforma de DeepBops, va a poder dejar documentado quién se logueó, qué acciones hizo, en qué momento, etcétera, etcétera. Entonces. Es un desafío, sí, de alguna manera seguir, conocer a los usuarios, pero estos usuarios no son inexpertos, no saben que están tocando, sino que son usuarios calificados y como usuarios calificados desde el punto de vista de la seguridad son más riesgosos también porque. Siempre está la curiosidad, no de este tipo de usuarios. Entonces, sus acciones son muy limitadas dentro de cada herramienta que deben usar. Sí, salvo que uno utiliza alguna herramienta open source que es totalmente abierta. Es muy raro en el mundo DevOps, pero la mayoría todas. Son capaces de segregar usuarios y trazar su actividad. Entonces los tenemos como bien acotados, no van a poder, por decir así, desintegrar. 23 herramientas, las configuraciones pasan por otro tipo de roles que no se les asignan a ellos.

    \item[\textbf{MUNTAABSKI FEDERICO:}] Bárbaro. Bueno, seguimos con la siguiente pregunta, que sería vos en tu rol de DevOps, vos, ¿qué tipo de incidentes de seguridad son los que más tienen presencia o que vos sentís que ¿Que aparecen más?

    \item[\textbf{PABLO:}] Claro, es una buena pregunta. En los proyectos donde estoy ahora son proyectos relativamente medianos y grandes, con lo cual ya hay un equipo más pequeñito que justamente son los ingenieros DevSecOps. Específicamente se dedican a la segurización de toda esta plataforma de herramientas integradas que les comento para que tengan una idea por cada proyecto IBM consensúa con su cliente. ¿Cuál va a ser el stack tecnológico? Es decir, hacen como una listita. De las herramientas que van a ser utilizadas en esa plataforma para sus proyectos DevOps en común. Entonces podemos tener cuatro o 5 clientes diferentes y tener plataformas diferentes, no con herramientas que se van integrando de manera diferente. entonces. Que. Hay una diversidad de herramientas y la segurización específica de estas plataformas está encargada a este grupo de DevSecOps. Ahora, lo que yo veo que habitualmente ocurre, o mejor dicho, se intenta de remediar Y probar antes de que ocurra son lo que se llaman denegaciones de servicio, porque IBM trabaja con varias nubes. IBM tiene su propia cloud, que es muy chiquitita, muy específica. IBM Cloud se llama, no es muy conocida, pero está abocada a ciertas empresas. Pero todas estas plataformas que les comento utilizan AWS, utilizan Google, utilizan Azure como proveedoras cloud más conocidas. Incluso hay un proyecto que tenemos donde están integradas las tres nubes, no todos los servicios de las tres nubes. Pero el cliente utiliza la base de datos en una nube, los servidores en otro, las APIs en otra. Bueno, entonces ese tipo de diagramas que está bueno, es interesante ver y les trae un poco de dolor de cabeza a los DevSecOps, a los ingenios DevSecOps. porque. Tienen muchos puntos de ataque y todos aquellos que son sí, sobre todo lo que se llaman. Api Gateway, no que por ejemplo en AWS son puntos donde se abre de alguna manera un puerto para que la información de esa infraestructura dentro de la nube salga por ese puerto, ¿no? Bueno, todo esto ese tráfico que sale y que entra necesita ser escaneado, analizado. Y también ese puerto que está abierto y que se usa para transaccionar. Muchas veces es una superficie de ataque donde desde afuera, bueno, se puede intentar alguna denegación de servicio, algún otro, alguna otra técnica. No, entonces es mi. En mi día a día lo que yo veo son más denegaciones de servicio. Directamente a esta plataforma pasó muy poco, pero sí a algún proveedor de esta plataforma. Nosotros usamos, como les digo, tres nubes, son proveedores externos a esta plataforma. alguno de los de las herramientas de esta plataforma están hosteadas en servidores on premis del cliente de IBM y otras son directamente plataformas como como les decía recién GitHub, GitLab. Entonces lo que puedo, lo que habitualmente ocurre no muy seguido, pero que por ejemplo se indisponibiliza la consola de GitLab y entonces no se pueden lanzar pipelines. Bueno, no es un problema directo de IBM y de sus herramientas, sino más bien que estamos consumiendo los servicios de GitLab y de pronto, bueno, indisponibilizaron la consola a través de denegación de servicios. No, eso sí ocurre, no muy seguido, pero ocurre. ¿Cómo es decir, como la pregunta iba a una técnica de? Vamos a decir así de seguridad y les pongo un ejemplo. A veces otras cosas que ocurren es que, por ejemplo, no sé, algún servicio de AWS está caído o la región donde tenemos los servicios, bueno, no está funcionando. Ok, eso no sería un ataque propiamente, pero bueno, tenemos inconvenientes para poder hacer funcionar toda la plataforma de DevOps, ¿no? Pero el ataque más conocido de negación de servicio.

    \item[\textbf{MUNTAABSKI FEDERICO:}] Claro. Perfecto, bueno, pasamos a la siguiente pregunta. Una vez que. que vos planteas esto de que usan un entorno, o sea en tu trabajo usas un entorno de múltiples nubes o que usas diferentes servicios, como mencionaste, Git DevOps o GitLab, no me acuerdo bien cual era...

    \item[\textbf{PABLO:}] Correcto, GitLab, sí.

    \item[\textbf{MUNTAABSKI FEDERICO:}] ¿Qué forma tienen ustedes de visibilidad del equipo de trabajo? O sea, ¿cómo ustedes aseguran de que el grupo de trabajo o los usuarios que están trabajando en los proyectos ¿Puedan ustedes controlar los procesos, qué herramientas, o sea, cómo se guían con eso?

    \item[\textbf{PABLO:}] Yeah. Bueno. Aquí nosotros usamos, mi equipo, los ingenieros DevOps, pero también hay un equipo en particular que es el equipo de observabilidad, así le llamamos. Tenemos dos tipos de herramientas, digamos, usamos por un lado, Hacemos lo que se llama una herramienta que usamos, por ejemplo, Grafana. Es una de las herramientas que utilizamos para crear dashboard a partir de datos que se censan en diferentes equipos de las diferentes herramientas que conforman la plataforma, no por ejemplo. Como les decía, GitLab, nosotros no lo tenemos instalado un premise dentro de uno de los clientes, sino que usamos directamente su servicio cloud. Entonces nosotros a través de Grafana miramos una como si fuera una pantalla de check que tiene GitLab para ver que sus servicios estén. Todos levantados, es como si fuera un HTML donde vos tenés en verde todos los servicios fundamentales de GitLab. Eso es lo que tienen habitualmente las plataformas cloud, que te muestran su estado de funcionamiento. Las nubes también lo tienen. Pero cuando usamos una herramienta que tenemos instalada en algún servidor on-premise, sea del cliente o sea de IBM. Nosotros tenemos un agente. Grafana en sí trabaja con otro tipo de bases de datos más este, digamos, no relacionales, como puede ser Elasticsearch o como puede ser Prometheus. Entonces, a esas bases donde se junta y se acopian datos, se le integran. Pequeños agentes, como podrían ser los agentes de Telegraph y esos pequeños agentes instalados en los servidores donde se realizan las acciones de los usuarios, nos ayudan a recolectar datos. Esos agentes, bueno. No son, no son agentes de IA agéntica, pero básicamente son como si fueran aplicaciones, muchas en Java, muchas en Python, que van a ir tomando los datos de diferentes variables de entorno de el sistema operativo donde vos lo instales, ¿no? Hay muchas plataformas de observabilidad que tienen este tipo de agéntica, porque no sé, la más emblemática, ya tiene como 30 años, es nagios. Entonces, vos gracias a ese agente que le reporta a tu servidor central, tenés los consumos de recursos, las transacciones que ocurren allí. puedes monitorear alguna carpeta, algún archivo específico en orden a esto que charlábamos recién, si hay algún cambio, alguna modificación. Bueno, volviendo a tu pregunta, nosotros lo que tenemos es una herramienta de observabilidad que en su primera capa, en la capa vamos a decir así gráfica, Es grafana. Esa es la herramienta que me que nos permite crear gráficos y después en su capa backend, bueno, vamos a tener una serie de servidores que van a recibir la información de diferentes agentes y esos agentes, luego de reportar la información, las vamos a guardar en una base. Habitualmente una base llamada Prometheus, que es una base no relacional que le permite a Grafana mirar esos datos y armar gráficos con esos datos. no? Entonces yo sé que hay 20 usuarios conectados en el servidor en el servidor A. porque hay un agente en el servidor que está reportando. Bueno, qué se está haciendo ahí, cuántos usuarios hay, qué cantidad de memoria de almacenamiento se está utilizando, qué archivos están escribiendo, etcétera. No son herramientas específicamente de seguridad, pero las usamos. Para hacer lo que se llama observabilidad, que justamente es eso, no observar qué es lo que se hace ahí, no y de qué manera, etcétera. Entonces nosotros usamos eso. Después el equipo de DevSecOps, es decir, los ingenieros de seguridad para este mismo proyecto, tienen otras herramientas un poco más sofisticadas como para ver si uno de esos usuarios no quiso acceder a una carpeta que no le correspondía, por ejemplo.

    \item[\textbf{MUNTAABSKI FEDERICO:}] Claro, entiendo.

    \item[\textbf{MOCIULSKY SANTIAGO BERNARDO:}] Ok te puedo hacer una pregunta con respecto a esto en lo que vos nombraste lo más en el entorno One premise lo de los múltiples agentes eso como output daría la información de la actividad que estuvieron haciendo on premise Pero cómo se procede después O sea cómo sería el paso a paso que se hace

    \item[\textbf{PABLO:}] sí. Sí. It.

    \item[\textbf{MOCIULSKY SANTIAGO BERNARDO:}] Después de que llega esa información, se actúa al momento o se actúa, cuánto se tarda en resolver ese problema.

    \item[\textbf{PABLO:}] Bien. Perfecto, es decir, una cosa es la recolección. De hecho, la observabilidad desde el punto de vista de infraestructura implica que vos puedas ver, ¿no? Entonces, esto que les cuento es, ni más ni menos, la posibilidad de poder ver en tiempo real lo que está sucediendo. Eso es lo que te permite Grafana con Prometheus y, por ejemplo, los agentes de Telegraph. Entonces vos tenés que esos agentes le reportan un server central, ese server central manda los datos a la base de datos que va creciendo porque es no relacional y Grafana, de acuerdo a diferentes vamos a decir query del tipo JSON, etcétera, vos podés armar gráficos con grafana. Entonces en tiempo real vos ves qué es lo que ocurre. Ahora, todo esto que te estoy contando no fue pensado originalmente para hacer seguridad, sino para hacer observabilidad. Así todo, con este tipo de Mirada observabilidad, nosotros podemos detectar que de pronto hay un usuario que está queriendo acceder a una carpeta que no le corresponde y entonces en mi dashboard de Grafana me empiezan a aparecer mucho registro o qué sé yo, no sé, como si yo les dijera permission de night. Sin muchos registros de tal usuario que tienen esta notificación en particular, no o k en ese caso nosotros. Y acá viene un poco la salvedad, reportaríamos el problema al equipo de seguridad. Nosotros como ingenieros DevOps en un proyecto de nuevo grande, como es este que, como son estos que les estoy contando, no accionamos directamente en lo que es, no sé, limitarle los Permisos a este usuario, porque nuestro scope de trabajo es la seguridad en la plataforma, es decir, si alguien, si un usuario me estuviera queriendo cambiar la configuración de integración entre 2 herramientas, ahí sí nosotros podemos de alguna manera desactivar ese usuario. Lo desactivamos de esa plataforma, sí le cortamos la sesión y que eso sería una acción de remediación inmediata más afín a mi rol. Ahora, si ese usuario está queriendo escribir una carpeta que no le corresponde, bueno, nosotros en este tipo de proyectos no entraríamos tanto en juego, pero sí podríamos reportarlo al equipo de DevSecOps de seguridad y ahí habitualmente, bueno, sí, hay algún tipo de aviso, generalmente Y. Si el usuario, ya les digo, son usuarios administradores, es decir, o son usuarios que tienen expertise técnico. Entonces, en la terminal donde está queriendo hacer esa acción, suele aparecerle algún tipo de cartel, no como para que sepa que está haciendo una acción que no está permitida y ellos tienen. Herramientas donde ese mismo usuario se va, se va copiando de ese usuario las diferentes acciones con prohibiciones. No, por ejemplo, quiere ejecutar, no sé, la sentencia sudo o el comando sudo en un terminal donde no está permitido. Bueno. Los ingenieros de DevSecOps tienen herramientas que le reportan a los gerentes directos de ese usuario que está queriendo hacer una acción ilegal todas esas acciones que efectivamente no debe poder hacer, no entonces. En este punto podría averiguarles si hay en tiempo real algo que ellos hagan en relación a un usuario que ya está en el sistema, distinto a que vos me digas, encontraron un proceso que nadie sabe de qué software es, eso es otro digamos ¿Otra acción, no? Yo sé que en este caso sí, los ingenieros de DevSecOps para este tipo de proyectos tienen la potestad de matar ese proceso, lo destruyen, hacen un kill de ese proceso y se acabó ese proceso. ¿Por qué? Porque no pertenece a un catálogo. No sabemos de dónde salió ese proceso, podría ser un malware. Pero referido a un usuario. Sé que le dan aviso, sé que le avisan al gerente directo, no sé si hacen algo más, yo entiendo que no, porque ese usuario está identificado, se conoce quién es, etcétera. Entonces, en tiempo real, volviendo a la pregunta de Santiago, en tiempo real, yo como ingeniero DevOps no hago nada, si no es básicamente una Sí, un ataque, por decir así, a la integración de 2 herramientas dentro de una misma plataforma. En el caso de que veamos otra cosa, se lo reportamos al equipo de seguridad que él tiene. Ellos tienen más herramientas, incluso para hacer escaneo de tráfico o escaneo de ciertas acciones. Dentro de un mismo ser.

    \item[\textbf{MUNTAABSKI FEDERICO:}] Claro, entiendo. Perfecto. Bueno, paso a la siguiente pregunta que es, o sea con lo que mencionaste, la observabilidad, o sea que tu área o el área en el que estás vos se dedica a lo que es O tienen herramientas que permiten observar el comportamiento, pero como que la parte de la seguridad la delegan a otro área.

    \item[\textbf{PABLO:}] Claro, la seguridad referida a la acción de los usuarios, la seguridad en cuanto a la integración de herramientas si la manejamos nosotros, pero ahí no hay otros usuarios que no seamos nosotros. Si yo tengo que unir Jenkins, que es una herramienta de despliegue con un ambiente de un servidor que está

    \item[\textbf{MUNTAABSKI FEDERICO:}] Alexander su usuario. Claro.

    \item[\textbf{PABLO:}] WS, yo hago la integración y hago que esa integración sea segura habilitándole, bueno, secrets, encriptación de datos en esa integración para que las 2 plataformas se vean entre sí. Eso lo hago yo ahora. Yo lo que cuido o mi equipo lo que cuida es la seguridad de la plataforma y de esas integraciones. Sí, ¿qué pasa dentro de cada una de esas plataformas? Ahí es donde de alguna manera tenemos divididos los roles con los ingenieros DevSecOps.

    \item[\textbf{MOCIULSKY SANTIAGO BERNARDO:}] Y entonces volviendo un poco desde que ustedes hacen una petición o sea ven algo y lo reportan a seguridad al área de ciberseguridad ¿Cuánto tarda en darle un feedback o solucionar ese problema que se está teniendo aproximadamente de tu experiencia?

    \item[\textbf{PABLO:}] Sí. Y buena pregunta. Mira, nosotros ahora en este proyecto que tengo en la mente que les estoy contando. Ya no tenemos software de tickets, sino que nos manejamos con una metodología Shyle muy parecida a Scrum, aunque no somos desarrolladores y utilizamos una herramienta que se llama Jira. Capaz que la conocen. Jira es como una herramienta donde vos. Las más diferentes peticiones. En forma de tarjetas y esas tarjetas van cambiando de estado en la medida en que los equipos de trabajo van tomando y trabajando sobre esos temas. Si yo te dijera un tiempo y ponele, calculale, voy a ser exagerado, pero calculale. En una, no sé, un tema que podría tener una severidad o no sé, una severidad media, quizás 2 días hábiles. Lo que tiene Gira como herramienta y en el cliente donde estoy trabajando, en el cliente de IBM, no está demasiado maduro esta parte, este procedimiento, con lo cual. A veces las tarjetas no siguen un tiempo determinado. No es que todas se cierran a los 3 días hábiles, 4 días hábiles. No está muy automatizado eso. Jira es muy flexible, entonces dependemos mucho de. Bueno, de que responda otro el otro equipo, que el otro equipo no tenga prioridades más altas, pero bueno, si quieren un número, pongámosle así, 2 días hábiles en recibir una respuesta del equipo de seguridad.

    \item[\textbf{MOCIULSKY SANTIAGO BERNARDO:}] Okay. Y no hace ninguna acción intermedia como para mitigar el problema. Volvamos un poco más orientado a nuestro problema. Estamos tratando de ver. Por ejemplo, un usuario que tiene permisos hace actividad sospechosa. Ustedes lo ven por el sistema de observabilidad. No hace ninguna acción inmediata o

    \item[\textbf{PABLO:}] Sí.

    \item[\textbf{MOCIULSKY SANTIAGO BERNARDO:}] Temprana para tratar de mitigar ese problema o tienen que esperar hasta esos 2 días.

    \item[\textbf{PABLO:}] Claro, nosotros como ingenieros DevOps no hacemos ningún inconveniente, ¿no? Lo que sí podemos hacer es crear una tarjeta sí con una severidad alta, reportando ese inconveniente en el servidor que lo vimos, pero tampoco nosotros somos el canal uno o la mesa de ayuda tecnológica, no , sino que somos un equipo.

    \item[\textbf{MOCIULSKY SANTIAGO BERNARDO:}] Hmm.

    \item[\textbf{PABLO:}] Uno de tantos de los que trabaja en el proyecto, que con una de sus herramientas que sirve para otra cosa, detectó un comportamiento anómalo. Está bien, ese sería como el escenario. Entonces, sí, si fuera una tarjeta de carácter urgente o crítica, quizás. En 2:03 horas, bueno, el equipo de seguridad se puso a mirar el tema. No creo, primero que nunca no pasó, pero no creo que sea mucho más rápido porque para que se den una idea, los ingenieros DevOps son los especialistas en seguridad. para los proyectos DevOps, pero no son los únicos que hacen seguridad para la compañía o para el cliente de IBM, sino que después están los equipos de seguridad estándar, seguridad de usuario, seguridad de red, ¿no? Entonces imagínense que todos esos están permanentemente mirando los comportamiento de toda la red sí del cliente Entonces lo más probable es que este tema de usuarios en particular salga o lo detecten bueno seguridad de usuarios y perfiles por ejemplo con sus propias herramientas que son herramientas específicas para ver qué hacen los usuarios en la red Sí, me explico. Es decir, el problema que estamos analizando puede darse en una infraestructura DevOps, pero luego se va a ir corriendo como problema a los equipos que naturalmente lo van a tener que resolver.

    \item[\textbf{MOCIULSKY SANTIAGO BERNARDO:}] Sí.

    \item[\textbf{PABLO:}] Que son los equipos tradicionales de seguridad.

    \item[\textbf{MOCIULSKY SANTIAGO BERNARDO:}] Okay, here you.

    \item[\textbf{PABLO:}] Por eso los tiempos que yo te puedo decir son, pero 3 horas para resolver un problema de seguridad, tienen un agujero por todos lados. Claro, porque son tiempos de problemas que no son los naturales a nuestro, básicamente a nuestro expertise, diferente a que yo tuviera un problema de... Integración de una herramienta, yo veo que integración de una herramienta directamente se perdió la integración o yo veo que la en la integración de una herramienta con una nube hay un montón de request permanentemente request. y entonces ahí sí nosotros salimos a investigar qué pasa porque no se está tirando ningún deploy a AWS pero de Jenkin tengo 50, 70 request por segundo, para algo pasa acá, ahí sí nos ponemos a analizar qué está sucediendo y bueno ahí puede haber, puede haber algún malware ¿Puede haber algún tipo de actividad anómala, pero ya no es un usuario, no?

    \item[\textbf{MOCIULSKY SANTIAGO BERNARDO:}] Yeah.

    \item[\textbf{MUNTAABSKI FEDERICO:}] Genial, bueno, ahora pasamos a otra pregunta.

    \item[\textbf{PABLO:}] Done.

    \item[\textbf{MUNTAABSKI FEDERICO:}] Vos, teniendo en cuenta las herramientas que mencionaste, Jenkins, Grafana, también esta metodología de trabajo con Jira y todo esto. Si vos pudieses, por ejemplo, proponer alguna mejora para este sistema o cómo lo cambiarías para que a vos se te facilite más el trabajo de, por lo menos la observabilidad de sentir que uno puede tomar más acciones o que se tiene que preocupar por menos cosas.

    \item[\textbf{PABLO:}] BM. Bien, buena pregunta. Sí, a ver. Acá la respuesta que te puedo dar es la siguiente. Dentro de lo que hace mi equipo, que como les digo, las los objetivos del ingeniero DevOps es sostener, es decir, darle mantenimiento a cada una de las herramientas del ciclo de vida DevOps, no herramientas de despliegue, herramientas de creación de imágenes, de testeo, etcétera. Ese es nuestro trabajo. Y la integración de esas herramientas para que todos trabajen en equipo. La seguridad de esas herramientas ya está como salvaguardada dentro de las acciones que realizamos, ¿no? Si yo tuviera que dar o pedir una mejora, podrían ser 2 caminos, o pido una mejora de una incluir una herramienta nueva. Porque vamos a tener nuevas, nuevas tareas relacionadas más con la seguridad, entonces ahí yo buscaría algún tipo de herramienta para poder ponerle escanear el tráfico en nuestra plataforma, eso podría ser una alternativa, podríamos nosotros hoy en día no tenemos permiso para usar ni nmap ni Wireshark Pero podría ser una alternativa, pero sería como adicionar a mi equipo tareas que actualmente no las tenemos definidas. Es decir, ¿quién escanea la red en este cliente donde trabajo? Bueno, la escanea seguridad de redes. Es decir, ellos son los que hacen esa tarea. Y. En mi equipo en particular, digamos, actualmente tenemos las herramientas adecuadas para hacer lo que está estipulado que nosotros hagamos. Yo lo que podría pedir sería una mejora de esas herramientas. No sé, en vez de usar, qué sé yo, Jenkins, que es una herramienta open source que no tenés una empresa que está detrás dándote soporte, quizás la cambiaría por no sé, GitHub. Action o GitLab, donde la plataforma y la empresa que está atrás es mucho más potente que el grupo de desarrolladores de Jenkins, no decir se nos desconfigura Jenkins y tenemos que salir a investigar a ver qué se tocó y cómo remediarlo, porque es open source, no, entonces yo haría eso. cambiaría a un, por supuesto, la empresa utiliza herramientas open source porque no tienen una licencia paga, se ahorran mucho dinero por eso, ¿no? Yo haría ese tipo de cambios. Cambiaría una herramienta por otra para mejorar nuestro trabajo. Desde el punto de vista de la seguridad, nosotros lo único que integramos es una herramienta que se llama Sonarcube. Es una herramienta específica para validar. Bueno, nos permite hacer ciertos testeos. Jenkin también nos permite hacer ciertos testeos de imágenes, etcétera. No hacemos escaneos de imágenes de containers, por ejemplo, pero bueno, es una herramienta que nos facilita los testeos. No, entonces ahí tenés otra otra acción que realiza mi equipo relacionada con la seguridad en un proyecto. Pero los cambios de herramientas serían esos. Esos serían las mejoras que podría utilizar. No saldría a pedir una herramienta específicamente de seguridad porque está fuera de nuestros COP. ¿Diferente a que yo trabajar en un equipo de DevSecOps y ahí mi ciento por ciento mi día sería hacer seguridad en muchas capas, no? En este caso no, yo hago algunas cosas de seguridad, algunas. Puedo conocer más, pero solo algunas por mi rol.

    \item[\textbf{MOCIULSKY SANTIAGO BERNARDO:}] Okay, thank you.

    \item[\textbf{PABLO:}] No sé si te respondí.

    \item[\textbf{MOCIULSKY SANTIAGO BERNARDO:}] sí se entendió perfecto.

    \item[\textbf{MUNTAABSKI FEDERICO:}] sí, entonces es perfecto.

    \item[\textbf{PABLO:}] Okay.

    \item[\textbf{MOCIULSKY SANTIAGO BERNARDO:}] Ahora te quería consultar más relacionado a ustedes en el acceso diario a los sistemas que tienen, cómo es el tema del control de acceso a identidades en su día a día trabajando en datos confidenciales, en la nube, aplicaciones que son críticas, etcétera.

    \item[\textbf{PABLO:}] Sí. Perfecto, y si eso que vos preguntas, Santiago, es re importante, es re importante tanto para mi empleador, IBM, como para su cliente. De hecho, antes de que comenzara este proyecto, esas 2 empresas firman un acuerdo y un contrato. En particular, donde los empleados de ambas compañías que van a. Estar trabajando dentro de este proyecto. Van a cumplir ciertos procesos para asegurar que la información no se filtre. Entonces, desde haber firmado un acuerdo de confidencialidad, es decir, yo tuve que firmar para este para este contrato puntualmente un acuerdo de confidencialidad donde. Mi empleador, pero de alguna manera mi persona se hace responsable de la información del cliente que pasa por mi computadora, por mis manos. Entonces yo estoy penalizado desde lo laboral o desde lo penal. Si de pronto filtro datos, no filtro datos del cliente, no en particular es un cliente financiero. Entonces ahí es como que está más. Más como auditado todo, ¿no? Entonces, primero eso, desde lo legal, todos los que formamos parte de este proyecto tenemos un acuerdo de confianza y confidencialidad con nuestro empleador y mi empleador con su cliente. Eso primera idea, con penalidades, etcétera. Segunda idea, la computadora que nosotros utilizamos. Sí, la que nos da, en este caso, el cliente. Si nosotros tenemos una de IBM, pero a partir de este proyecto tenemos una computadora específica para trabajar para este cliente. Esa computadora es lo que a mí me gusta decirle, es un dispositivo totalmente bobo, bobo en el sentido que no le puedes conectar nada. De hecho, si yo se las muestro, es una Lenovo, está buenísima y tiene mil puertos y no se conecta a ningún puerto, están todos los puertos limitados y su conexión solamente se da a través de esto. Esto que ustedes ven aquí es un dispositivo 3G, tiene un chip adentro, como si fuera un USB Y adentro tiene un chip Movistar y entonces la computadora tiene acceso a la red ni Internet a la red a través de este dispositivo. Es un módem. Sin este dispositivo la computadora tiene una imagen en su sistema operativo que no te permite usar ningún puerto, ni tarjeta de red ethernet, ni tarjeta Wi-Fi, no te funciona nada, vos prendes todo y no detecta nada tiene todo filtrado salvo cuando le pones este este módem 3g y ahí directamente se conecta vía 3g a través de un chip Movistar a la red del banco y entonces ahí empieza a pedir información usuario password y token del Cliente de este cliente IBM, que es un banco, es lo primero que te pide. Una vez que te pidió de eso, encripta la levanta una VPN, encripta la comunicación entre este equipo y los servidores de. Del banco y ahí recién, digamos, está segura la conexión entre mi dispositivo, esté donde esté en el globo terráqueo y los servidores del banco. Y ahí, si yo quiero acceder a mis aplicaciones, porque la máquina es de nuevo, es totalmente boba, no tiene instalado ningún. Más que creo que tiene Windows 11, más que Windows no tiene nada, no tiene un notepad, nada. No podés, No podés navegar, no podés escuchar música, no podés hacer nada. Una vez que está levantada la VPN se te habilitan tres íconos donde vos podés acceder a un escritorio remoto del tipo Citrix o BDI, digamos son escritorios remotos que se ejecutan no localmente sino en servidores de la compañía, es decir, una vez que la VPN está Activa, vos podés abrir un escritorio remoto y entonces esta máquina termina siendo como una terminal. que todos los recursos tiene como, no sé, 16 gigas, no sé, i7, una bestia, no puedes usar nada de eso y solo sirve para levantar una VPN y un escritorio remoto, incluso los datos que yo utilizo no están en el equipo, están en los servidores de la compañía del otro lado de la VPN. No, entonces, desde el punto de vista de la seguridad, eso es tremendo. Cada vez que yo tengo que trabajar, incluso cuando hago guardias por la noche, tengo que levantar todo este procedimiento hasta llegar al escritorio remoto que me permite conectarme con los servidores. Dentro del cliente, no, eso tarda. Y sí, tardó como 15 minutos en levantar la máquina, con suerte, ¿no? Como es todo módem 3G, hay lugares donde no tenés señal, entonces yo vivo en un edificio, con lo cual eso influye en las comunicaciones también, tenés que estar al lado de la ventana, etcétera. Y una vez que hago todo el caminito y accedo al escritorio remoto, Del cliente, recién ahí puedo acceder a las consolas de las herramientas que tengo que validar, configurar, etcétera. No, entonces sí, es muy fuerte el proceso de bueno de de logueo y autenticación, autenticación y autorización del usuario. Sin esta computadora, con estas configuraciones, no hay forma de entrar en este cliente, no tenés ni chance. ejaz.

    \item[\textbf{MUNTAABSKI FEDERICO:}] Perfecto.

    \item[\textbf{MOCIULSKY SANTIAGO BERNARDO:}] Genial y con respecto al tema de gestión de identidades un poco ya por curiosidad eh ustedes se manejan en la empresa con un yam o cómo manejan el tema de que verdaderamente sos vos y qué tanto les permite o controla esa gestión de identidades

    \item[\textbf{PABLO:}] Bien, buena pregunta. Mira, en este proyecto, es decir, una cosa es IBM, cuando yo trabajo para IBM, uso una computadora, la tengo de este lado. Cuando trabajo para el cliente IBM, uso la otra computadora, e incluso a veces uso las 2, porque Los chats son diferentes y yo tengo que hablar con mi jefe de IBM, con mi gerente, hablo por un chat. Si tengo que hablar con el product owner del proyecto donde trabajo, hablo con otro chat. Desde esa disociación de equipos, incluso estas empresas hacen eso para que los datos del banco no pasen por los dispositivo de IBM, es decir, para eso es que a los empleados nos dan dos dispositivos. Por supuesto es mucho más, vamos a decir agresivo la validación del lado de la computadora del banco, mucho más agresivo. Y ahí el banco hasta lo que yo sé utiliza un servicio tipo IAM de la nube de Azure o alguna integración con la nube de Azure, ellos se maneja mucho con productos Microsoft y la primera validación es local, es decir, la computadora cuando hago el login, cuando levanto el sistema operativo me pide usuario, password, y token, un token que es básicamente una aplicación que el banco me hizo instalar, configuró. En este banco en particular fue una aplicación que me hicieron bajar a mi celular personal, pero hace unos meses, el año pasado, trabajé para otro cliente de IBM, que era otro banco, e incluso me daban un celular. Era un Samsung A24 creo. ¿Para qué? Para que las aplicaciones móviles no pasaran por mis dispositivos personales, pasaran por los dispositivos del banco. Entonces yo venía con una computadora, un smartphone del banco, la computadora de IBM, entonces irse a un lado, el otro era un lío. ¿Este a ese extremo, no? Entonces, actualmente sí, hay una autorización local de la computadora con estos 3 elementos que te digo, usuario, password, token. Y una vez que se hace esa primera validación, ahí está la segunda validación con casi los mismos datos, pero para acceder a la VPN. Hay también en la máquina del banco. hay. ¿Cómo? Tercer factor de autenticación, datos biométricos, es decir, el face ID y la huella son elementos que esta Lenovo también tiene configurado por default. Bueno, y una vez que levantó la VPN, ya ahí, bueno. La tercera validación seguiría en las consolas de las herramientas donde tengo que ingresar. Por ejemplo, si yo tengo que ingresar a GitLab del banco, me va a volver a pedir un usuario y un password específicamente para entrar a los proyectos del banco. Así que básicamente tengo 3 validaciones de autenticación. Vos me dirás, bueno, pero son los mismos datos. Sí, son los mismos datos. Por ahí a veces cambian el factor de autenticación. Algunos son biométricos, otros son un código. En una aplicación específica. Pero sí, son por lo menos tres seguro. ¿Contra qué lo hace? Yo entiendo que lo hacen contra un servicio de IAM, de Azure. No estoy muy empapado, te lo podría averiguar, pero me parece que sí. Sobre todo para poder levantar todo lo que es la VPN particular del equipo. No. Yeah.

    \item[\textbf{MOCIULSKY SANTIAGO BERNARDO:}] Okay y como última consulta sobre este tema y al momento de que vos haces esa triple verificación tenés alguna limitación o algún control con respecto a la actividad que haces o a qué archivos accedes etcétera se tiene alguna revisión con respecto a eso

    \item[\textbf{PABLO:}] Bien, esa también es una buena pregunta porque mi proyecto, el proyecto de este banco, según la cultura DevOps, yo soy un grupo, estoy dentro de un grupo de ese proyecto, no es el único proyecto de todo el banco y es un proyecto Que está dentro de la órbita de la infraestructura del banco, donde hay otras personas, otros otros procesos, otras otras metodologías. El banco utiliza también ITIL, utiliza también Kanban. Es decir, hay otras metodologías, hay otra. Bueno, Scrum, hay otras metodologías en proyectos IT dentro de este banco, no, por ejemplo. Hace unos años yo trabajaba como sysadmin de Linux, ¿no? Antes de ser ingeniero de DevOps, hoy como ingeniero DevOps en este banco puedo ingresar algunos algunos servidores Linux. Sí, algunos son premis, algunos cloud puedo ingresar y yo tengo limitado para mi usuario. Cuando me hicieron el pasaje de grupo, me limitaron, los me cambiaron de perfil, entonces me me cambiaron los permisos. Entonces yo puedo ejecutar ciertos comandos sudo, es decir, es decir, pidiendo ser admin del equipo, pero solo ciertos comandos. Ya no soy más root de los equipos Linux que forman parte de mi proyecto, ¿no? Entonces, bueno, ahí ciertos cambios de variable entorno que no puedo hacer. Sé cómo hacerlo, sé cómo hacerlo, pero si yo intento hacerlo en la consola de tal server, bueno, se activa todo lo que les contaba un informe automático que le llega a mi gerente de IBM diciendo esta persona intentó acceder y ejecutar tal cosa a todo el usuario, eso no está permitido. La primera vez llega, la segunda se lo reportan a mi empleador y al cliente, digamos al product owner del cliente, y la tercera vez que lo hago seguramente me van a sancionar, me van a llamar la atención. Es decir, bueno, tres veces quisiste ser root y sabés que no es porque querés ser. si ahí empiezan a aparecer las alertas no Entonces sí mi cuando me cambiaron de grupo para este mismo proyecto me limitaron los permisos sobre todo en la administración de los servidores que venía dándole servicio No sí porque ahora yo ya no administro servers Yo administro aplicaciones, sí, o estas aplicaciones DevOps dentro de alguno de los servers.

    \item[\textbf{MOCIULSKY SANTIAGO BERNARDO:}] Ok, pero entonces estas acciones que se realizaban con respecto a seguridad se hacían al momento de que vos hacías la acción no tenían una revisión previa para confirmar de que vas por un lado que no tendrías que ir, por ejemplo.

    \item[\textbf{PABLO:}] Claro, no, en este caso que te digo, que es la ejecución de un comando dentro de un Linux, se hacían en tiempo real con herramientas que tenía la gente de seguridad de sistemas operativos, que no tienen nada que ver con seguridad de DevSecOps, es algo más amplio.

    \item[\textbf{MOCIULSKY SANTIAGO BERNARDO:}] Okay, genial.

    \item[\textbf{PABLO:}] A ellos le sale una alerta y esa alerta, es che, mira, tal usuario está haciendo esto, sí.

    \item[\textbf{MOCIULSKY SANTIAGO BERNARDO:}] Y con respecto a otros acciones que hacías, había algún tipo de revisión de las acciones o.

    \item[\textbf{PABLO:}] que no sea ejecución en un server, otro tipo de acciones. Bueno, por ejemplo, si en algún equipo Windows yo quisiera copiar carpetas o acceder a un disco, un disco compartido, está dentro de la red, pero compartido, que no tengo permisos, ahí también. Que sale el cartelito de permission the night. Es en tiempo real la validación, pero eso es una política que tiene ese servidor o esa carpeta y también la reiterada acción de esa situación se la reportan a mis jefes. Sí.

    \item[\textbf{MOCIULSKY SANTIAGO BERNARDO:}] Genial.

    \item[\textbf{MUNTAABSKI FEDERICO:}] Bueno, perfecto. Vamos a ir medio cerrando porque ya pasó bastante tiempo. Las últimas preguntas son más que nada de todo lo que hablamos. Si hay alguna pregunta que pensás que te deberíamos haber hecho y no te hicimos algo que se te ocurra ahora en el momento.

    \item[\textbf{PABLO:}] The win.

    \item[\textbf{MUNTAABSKI FEDERICO:}] Y nada, si querés agregar algo, agregar algo antes de finalizar.

    \item[\textbf{PABLO:}] Buena pregunta. Como alguna pregunta que me debieran haber hecho, no, la verdad que no se me ocurre. Fue bastante como amplia las preguntas que me hicieron y fuimos por diferentes partes de la seguridad de un proyecto. En particular, bueno, este que les contaba. del mundo DevOps simplemente agregar quizás o agregarles que actualmente en este banco nosotros no estamos integrando herramientas que tengan Inteligencia artificial porque hay una regulación del Banco Central de la República Argentina que limita el uso de la guía en Entidades financieras, entonces el Banco Central de la República Argentina, por ahora, le dice a los bancos que pueden usar y que no dentro de las nuevas tecnologías. Por ejemplo, un banco no puede tener bases de datos en nubes cloud, no hay, no está permitido, entonces por eso los bancos. que son bancos en Argentina y que están digamos categorizados así, tienen que cumplir esas disposiciones distintas, si yo le dijera Nubank o Brubank que son entidades financieras que pueden trabajar en Argentina pero no son Definidos como banco, no, entonces esas siguen otro tipo de, vamos a decir así, de reglas de juego. No, que ustedes sepan que la aplicación de inteligencia artificial en las entidades bancarias está por ahora restringido. Eso capaz que cambia, pero por ahora está restringido. Porque todavía no se sabe a ciencia cierta qué hace esas herramientas de IA con la información que tienen los bancos, que en el 99.9\% es información sensible. Sí, como no sabemos o que o como el mercado todavía no sabe. Está donde la herramienta de IA puede llegar o puede transaccionar si se la deja libre, por decir así. Entonces, por ahora no hay implementado software de IA, por ejemplo, en los proyectos donde yo trabajo, por ahora no hay, incluso tampoco para hacer monitoreo ni observabilidad. Pero eso es como anecdótico para que si ustedes siguen pensando este proyecto que ustedes tienen, bueno, quizás una limitante sea, bueno, el rubro financiero en Argentina no se lo van a poder vender, por ejemplo.

    \item[\textbf{MUNTAABSKI FEDERICO:}] Sí, esa es muy buena información. No estábamos al tanto de que el rubro financiero estaba regulado por el Banco Central, habías dicho, ¿verdad?

    \item[\textbf{PABLO:}] Sí, eso es cambiante, habría que averiguar, pero entiendo que no está permitido todavía.

    \item[\textbf{MUNTAABSKI FEDERICO:}] Perfecto, bueno, eso sería todo.
\end{description}
